% Appendix (O-RAN / BMPP-DQN track).
%
% \appendix switches \chapter numbering to letters (report class), so the
% single \chapter{} call below typesets as "Appendix A". Per-seed figures
% in Sections A.2/A.3 are read directly from each run's own summary.json
% (data/results_oran/, data/results_oran_rq3_single_timescale/) and are
% exactly the per-seed values the mean/CI figures in
% Chapter~\ref{chap:oran-results}'s tables aggregate -- included here, not
% in the main chapter, per this thesis's own Reproducibility quality gate
% (raw per-seed numbers, not only the aggregated summary).

\appendix

\chapter{Supplementary Material}
\label{chap:oran-appendix}

\section{Default Configuration Listing}
\label{sec:oran-appendix-config}

The full contents of \texttt{config/oran\_default.yaml}, the single
configuration file governing every result in
Chapter~\ref{chap:oran-results} (canonical $n=3$ run) unless stated
otherwise (the $n=10$ validation and the RQ3/sensitivity ablations each
change only the one parameter named where they are introduced). Every
constant's needs-validation status is disclosed inline, in this same file,
at the point each constant is defined -- not repeated separately here.

{\footnotesize
\begin{verbatim}
# Default Configuration for O-RAN / BMPP-DQN Experiments
# ORAN_BMPP_DQN_Concept_Note_v1.md -- a separate, additive research track
# alongside config/default.yaml's C-RAN track. No keys/sections here are
# read by cran_env/agents/training/evaluation, and vice versa.

# ============================================================
# Network Topology (Concept Note Section 10.3: single-gNB, focused scope;
# 2026-08-30 literature check of DQRL/OREO's own scenario scales -- see
# Section 10.3's own note -- shows this UE:RU ratio sits within
# precedented ranges, though the absolute RU count is still smaller than
# either paper's own scenario)
# ============================================================
network:
  n_ru: 4
  n_du: 1
  n_cu: 1
  n_ue: 8
  n_splits: 3           # 3GPP TR 38.801 Options 2/6/8, mapped to centralization
                          # levels c=0/1/2 -- Concept Note Section 10.2 (needs validation)
  area_size_m: 500
  carrier_freq_ghz: 3.5
  # 2026-09-24: raised from 20 -> 100 MHz. Empirical check (best-case
  # all-RU/max-power/equal-PRB policy, 300 steps): at 20 MHz, mean
  # aggregate demand (~101 Mbps, driven by this scenario's Poisson
  # burstiness) already approached mean achievable throughput (~77 Mbps)
  # even under that generous reference policy, leaving essentially no
  # headroom for any policy -- including the intended learned ones -- to
  # satisfy QoS. 100 MHz raises mean achievable throughput to ~244 Mbps
  # under the same reference policy (best-case per-UE satisfaction
  # 74%->81%), giving real headroom while still being a standard,
  # precedented channel bandwidth for this model's own n78-band carrier
  # frequency (3.5 GHz) -- not an arbitrary tuning choice. See
  # docs/oran_thesis_guide.md's QoS-metric note for the accompanying
  # per-UE-vs-all-UE metric distinction this same investigation surfaced.
  # 2026-10-01 re-check after fixing oran_env.py's PRB-sharing bug (each
  # co-served UE was previously scored against its serving RU's *full*
  # share, not divided among every UE that RU simultaneously served --
  # see Chapter 4's bug-fix disclosure): the ~384 Mbps/84% figures above
  # were computed under that inflated formula and are superseded by the
  # ~244 Mbps/81% figures now stated; 100 MHz still clears this headroom
  # check under the corrected formula (~2.4x margin), so the constant
  # itself did not need to change, only this comment's own numbers.
  bandwidth_mhz: 100
  noise_power_dbm: -102

# ============================================================
# Channel Model (simplified: log-distance path loss + Rayleigh fading only,
# no shadowing/temporal correlation -- Concept Note Section 5.1)
# ============================================================
channel:
  path_loss_exponent: 3.5

# ============================================================
# Traffic Model (trapezoidal daily envelope x Poisson per-UE arrivals,
# Concept Note Section 5.1; needs-validation placeholders per Section 10.6.
# 2026-08-30: lambda_peak/packet_size_bits below are now derived directly
# from 3GPP TR 38.864 Annex A's FTP Model 3 (0.5 MB packet size, 200 ms
# mean inter-arrival time -- a real, standard 3GPP Poisson traffic model),
# not a guess -- see Section 10.6's own note. floor_ratio and t1-t4 remain
# unvalidated placeholders: TR 38.864's own load scenarios (Table A-1) are
# instantaneous PRB-utilization snapshots, not tied to time-of-day, and its
# own study scope stops at "medium load" (<=50%), so it gives no floor:peak
# ratio or diurnal timing to derive these from.
# ============================================================
traffic:
  # 3GPP TR 38.864 Annex A FTP Model 3: 200ms mean inter-arrival = 5
  # arrivals/s = 0.5 per 0.1s step (peak Poisson rate, per UE)
  lambda_peak: 0.5
  floor_ratio: 0.2           # off-peak floor as a fraction of lambda_peak
                             # -- still an unvalidated placeholder (see note above)
  # 3GPP TR 38.864 Annex A FTP Model 3: 0.5 MB packet/file size = 4e6 bits
  packet_size_bits: 4.0e+6
  step_duration_s: 0.1       # matches the 10-100ms Near-RT RIC decision rate
  t1: 7.0                    # rise start (hour of day) -- still unvalidated
  t2: 10.0                   # plateau start -- still unvalidated
  t3: 20.0                   # plateau end -- still unvalidated
  t4: 23.0                   # fall end -- still unvalidated

# ============================================================
# Power Model (RU/DU/CU/fronthaul, monotonic in split centralization level;
# Concept Note Section 10.5). Most constants below remain needs-validation
# literature-style placeholders; p_max_dbm (below) is the one exception,
# now derived from a real small-cell O-RU datasheet -- see its own comment.
# ============================================================
power:
  ru:
    p_sleep_w: 2.0
    p_switch_w: 2.0           # RU activation-flip switching cost
    p_switch_split_w: 1.0     # RU split-change switching cost
    pa_efficiency: 0.25
    # 2026-08-30: 33 dBm = 2 W, matching the Benetel RAN550's real max TX
    # output power (total EIRP) for a Split-7.2x indoor small-cell O-RU
    # (was 30 dBm/1 W, an unvalidated placeholder) -- see Concept Note
    # Section 10.5's own note. Not a guess: same physical quantity, same
    # units, a real commercial small-cell O-RU datasheet.
    p_max_dbm: 33
    # Active-RU processing power per split centralization level c=0/1/2
    # (Concept Note Section 10.2), c=0 (Option 2) = most RU-side processing
    # down to c=2 (Option 8) = least. Previously not exposed here at all --
    # oran_env/oran_env.py never passed this array to ORANPowerModel, so
    # only its Python-side default (unchanged below) was ever actually
    # used; now genuinely config-driven, matching every other constructor
    # argument's existing convention.
    p_proc_by_split_w: [10.0, 6.0, 3.0]

  du:
    p_static_w: 50.0
    # DU processing power contributed per active RU at split level c
    # (increasing with c, since less-centralized splits push more
    # processing to the DU) -- same previously-not-wired gap as
    # ru.p_proc_by_split_w above.
    p_per_ru_by_split_w: [5.0, 10.0, 20.0]

  cu:
    p_static_w: 30.0
    p_dyn_per_ru_w: 1.0

  fronthaul:
    p_common_w: 10.0
    # Fronthaul power contributed per active RU at split level c
    # (increasing with c, since less-centralized splits need more
    # fronthaul bandwidth) -- same previously-not-wired gap as above.
    p_per_ru_by_split_w: [3.0, 6.0, 15.0]

# ============================================================
# DRL Algorithm (BMPP-DQN: two-timescale upper/lower Q-networks, no TD3 --
# Concept Note Section 5.2/10.4)
# ============================================================
algorithm:
  name: "bmpp_dqn"
  # 2026-09-29: reverted back to 500 (was briefly raised to 1000 on
  # 2026-09-26; see docs/daily_log.md's 2026-09-26 entry for that
  # rationale, still valid as a finding, not retracted) -- reverted at the
  # user's explicit request, back to this thesis's originally-reported
  # 3-seed/500-episode scale (Concept Note Section 6.5). The 5-seed/
  # 1000-episode expanded results (Section 6.6-6.7) remain archived at
  # data/results_oran_archive/5seed_1000ep_pre_exploration_fix_20260929/,
  # not discarded.
  max_episodes: 500
  max_steps_per_episode: 100

  # Two-timescale cadence: discrete (upper) decisions held for N lower-level
  # steps -- Concept Note Section 5.2 ("1-10s" upper vs "10-100ms" lower;
  # at step_duration_s=0.1, N=10 is the low end of that range, needs validation)
  upper_level_period_steps: 10

  # Replay buffers: separate lengths per Concept Note Section 5.2
  # ("longer replay buffers" for upper-level, "faster updates" for lower-level)
  upper_buffer_size: 200000
  lower_buffer_size: 1000000
  batch_size: 128
  min_buffer_size: 2000

  lr_discrete: 1.0e-4
  lr_actor: 3.0e-4
  lr_critic: 3.0e-4        # Used by the DDPG baseline (oran_agents/ddpg_agent.py)

  # Smaller than the C-RAN track's [256, 128] -- matches this track's
  # smaller single-gNB scope (Concept Note Section 5.1)
  hidden_dims: [128, 64]
  activation: "relu"
  use_layer_norm: true

  gamma: 0.99
  tau: 0.005

  epsilon_start: 1.0
  epsilon_end: 0.01
  epsilon_decay: 0.995

  continuous_noise_std: 0.1
  continuous_noise_std_end: 0.01

  gradient_clip_norm: 1.0

  # Tractability cap for the flat MP-DQN baseline's joint action space
  # (2^n_ru * n_splits^n_ru); mirrors config/default.yaml's analogous cap
  # for its own MP-DQN baseline.
  max_n_ru_for_flat_joint_action: 6

# ============================================================
# Reward Function Weights (needs validation -- see Concept Note
# Section 10.7; unlike every power/traffic/split constant in this file,
# these three weights have never been subjected to a literature check or
# a sensitivity analysis, despite directly defining the optimization
# objective every result in this thesis is measured against.
# oran_env/oran_env.py's step(): reward = alpha_energy * (throughput_mbps
# / total_power_w) - beta_qos * (qos_shortfall_mbps) - gamma_switch *
# (switching_events). Chosen only to preserve a reasonable relative
# magnitude -- beta_qos=10 so a QoS violation (Mbps of shortfall)
# dominates the throughput-per-watt term, gamma_switch=0.5 as a mild,
# not dominant, switching penalty -- not derived from any source.
# 2026-09 sensitivity sweep (Concept Note Section 6.9) perturbs the
# power-model constants these weights are combined with, not the
# weights themselves; the weights remain open for future work.
# ============================================================
reward:
  alpha_energy: 1.0
  beta_qos: 10.0
  gamma_switch: 0.5

# ============================================================
# Evaluation (Concept Note Section 5.3: 3 seeds, per RQ2)
# ============================================================
evaluation:
  eval_freq: 50
  n_eval_episodes: 5
  n_random_seeds: 3   # [42, 123, 456] -- Concept Note Section 5.3 "3 random seeds"
  save_checkpoints: true
  checkpoint_freq: 250

# ============================================================
# Hardware
# ============================================================
hardware:
  device: "cuda"   # cuda or cpu (falls back to cpu automatically if no GPU
                    # is available; oran_agents/*.py read this as each
                    # agent's device default)
\end{verbatim}
}

\section{Per-Seed Raw Results: Canonical Comparison (\texorpdfstring{$n=3$}{n=3})}
\label{sec:oran-appendix-perseed-canonical}

Table~\ref{tab:oran-appendix-canonical} gives the individual per-seed final
evaluation metrics each method's mean row in
Table~\ref{tab:oran-convergence} (Section~\ref{sec:oran-performance-comparison})
aggregates, read directly from each run's own
\texttt{summary.json} (\texttt{data/results\_oran/}). ``QoS Rate'' and
``Per-UE QoS'' follow the same two definitions given in
Section~\ref{sec:oran-experimental-setup}.

\begin{table}[htbp]
\centering
\caption{Per-seed final evaluation metrics, canonical $n=3$ run, all 4
  methods.}
\label{tab:oran-appendix-canonical}
\resizebox{\textwidth}{!}{%
\begin{tabular}{llrrrrrr}
\toprule
Algorithm & Seed & Reward & Power (W) & QoS Rate (\%) & Per-UE QoS (\%) & Switching & Throughput (Mbps) \\
\midrule
BMPP-DQN & 42  & $-$69424.60 & 151.49 & 22.8 & 78.43 & 0.236 & 230.37 \\
BMPP-DQN & 123 & $-$67526.27 & 145.08 & 22.0 & 77.72 & 0.184 & 194.92 \\
BMPP-DQN & 456 & $-$72152.45 & 131.35 & 22.2 & 77.53 & 0.124 & 178.55 \\
\addlinespace
DDPG & 42  & $-$87914.28 & 207.14 & 21.8 & 76.78 & 0.040 & 247.24 \\
DDPG & 123 & $-$100234.78 & 206.07 & 21.0 & 74.23 & 0.040 & 197.37 \\
DDPG & 456 & $-$88119.47 & 204.20 & 20.4 & 76.70 & 0.040 & 249.02 \\
\addlinespace
DQN & 42  & $-$56464.21 & 234.47 & 23.6 & 79.00 & 0.654 & 260.76 \\
DQN & 123 & $-$58656.24 & 196.20 & 22.4 & 79.00 & 1.128 & 251.89 \\
DQN & 456 & $-$58649.02 & 203.00 & 22.4 & 79.00 & 0.856 & 246.02 \\
\addlinespace
MP-DQN & 42  & $-$101610.51 & 142.30 & 20.0 & 73.00 & 0.468 & 52.33 \\
MP-DQN & 123 & $-$77158.25 & 158.76 & 22.2 & 75.33 & 0.354 & 122.69 \\
MP-DQN & 456 & $-$67004.92 & 149.41 & 28.4 & 78.83 & 0.040 & 260.51 \\
\bottomrule
\end{tabular}%
}
\end{table}

{\sloppy
DDPG's identical $0.040$ switching frequency across all three seeds is
a clean architectural fact, not an artifact of any one seed: unlike an
earlier version of this run (made under an environment-physics bug since
fixed, see Chapter~\ref{chap:oran-results}'s bug-fix disclosure), no seed
here shows a collapsed, nonsensically-low throughput alongside DDPG's
uniform switching count -- all three seeds' throughput (197-249 Mbps) and
power (204-207 W) are mutually consistent. The $0.040$ figure is simply
what Section~\ref{sec:oran-performance-comparison} already attributes it
to: DDPG's purely continuous relaxation architecturally never flips a
discrete decision in the sense \texttt{switching\_events} counts.
MP-DQN's seed 456 shows the same $0.040$ value by coincidence of its own
learned policy settling on a stable split/activation choice for that
seed specifically, not for the same architectural reason -- MP-DQN can
and does switch (seeds 42/123 show $0.468$/$0.354$).
\par}

\section{Per-Seed Raw Results: RQ3 Two-Timescale Ablation}
\label{sec:oran-appendix-perseed-rq3}

{\sloppy
Table~\ref{tab:oran-appendix-rq3} gives the per-seed figures underlying
Section~\ref{sec:oran-rq3}'s RQ3 ablation, for the same three seeds under
both configurations: the two-timescale default
(\texttt{upper\_level\_period\_steps=10}, identical to the BMPP-DQN rows
of Table~\ref{tab:oran-appendix-canonical}) and the single-timescale
ablation (\texttt{upper\_level\_period\_steps=1},
\texttt{config/\allowbreak rq3\_single\_timescale.yaml}).
\par}

\begin{table}[htbp]
\centering
\caption{Per-seed final evaluation metrics, RQ3 ablation: two-timescale
  (proposed, $N=10$) vs.\ single-timescale ($N=1$) BMPP-DQN.}
\label{tab:oran-appendix-rq3}
\resizebox{\textwidth}{!}{%
\begin{tabular}{llrrrrrr}
\toprule
Configuration & Seed & Reward & Power (W) & QoS Rate (\%) & Per-UE QoS (\%) & Switching & Throughput (Mbps) \\
\midrule
Two-Timescale ($N{=}10$) & 42  & $-$69424.60 & 151.49 & 22.8 & 78.43 & 0.236 & 230.37 \\
Two-Timescale ($N{=}10$) & 123 & $-$67526.27 & 145.08 & 22.0 & 77.72 & 0.184 & 194.92 \\
Two-Timescale ($N{=}10$) & 456 & $-$72152.45 & 131.35 & 22.2 & 77.53 & 0.124 & 178.55 \\
\addlinespace
Single-Timescale ($N{=}1$) & 42  & $-$64889.10 & 233.19 & 23.8 & 79.00 & 0.566 & 275.84 \\
Single-Timescale ($N{=}1$) & 123 & $-$61302.61 & 239.14 & 23.2 & 77.68 & 0.846 & 233.94 \\
Single-Timescale ($N{=}1$) & 456 & $-$59417.14 & 232.37 & 20.6 & 78.65 & 0.520 & 231.03 \\
\bottomrule
\end{tabular}%
}
\end{table}

Every single-timescale seed's switching frequency exceeds every
two-timescale seed's, with no overlap (minimum single-timescale value
$0.520$ vs.\ maximum two-timescale value $0.236$) -- the clearest
seed-level evidence for Section~\ref{sec:oran-rq3}'s conclusion that the
two-timescale design's effect on switching stability is consistent across
seeds individually, not only in the three-seed mean. Unlike the earlier,
pre-bug-fix version of this ablation, the two configurations here are
\emph{not} reward/power/QoS-neutral: every single-timescale seed also
shows higher power and (mostly) higher reward than every two-timescale
seed, a genuine trade-off discussed in Section~\ref{sec:oran-rq3}, not a
free reduction in switching.

\section{Per-Seed Raw Results: Fair-Cadence Control}
\label{sec:oran-appendix-perseed-cadence}

Table~\ref{tab:oran-appendix-cadence} gives the per-seed figures
underlying Section~\ref{sec:oran-cadence-control}'s fair-cadence control:
DQN and MP-DQN retrained and re-evaluated with the same
\texttt{upper\_level\_period\_steps=10} discrete-decision hold BMPP-DQN
uses by default (\texttt{discrete\_hold\_steps=10}), rather than
re-deciding \texttt{ru\_on}/\texttt{split} every step as in
Table~\ref{tab:oran-appendix-canonical}'s canonical rows for these same
two methods.

\begin{table}[htbp]
\centering
\caption{Per-seed final evaluation metrics, fair-cadence control: DQN and
  MP-DQN retrained with BMPP-DQN's own $N=10$ discrete hold.}
\label{tab:oran-appendix-cadence}
\resizebox{\textwidth}{!}{%
\begin{tabular}{llrrrrrr}
\toprule
Algorithm ($N{=}10$) & Seed & Reward & Power (W) & QoS Rate (\%) & Per-UE QoS (\%) & Switching & Throughput (Mbps) \\
\midrule
DQN & 42  & $-$104047.70 & 104.30 & 19.4 & 72.90 & 0.096 & 39.50 \\
DQN & 123 & $-$61684.90 & 148.00 & 23.6 & 77.80 & 0.084 & 224.40 \\
DQN & 456 & $-$64477.30 & 123.60 & 23.2 & 78.30 & 0.054 & 212.50 \\
\addlinespace
MP-DQN & 42  & $-$64403.50 & 141.00 & 23.8 & 78.40 & 0.030 & 220.40 \\
MP-DQN & 123 & $-$106382.30 & 132.40 & 18.6 & 72.70 & 0.086 & 2.30 \\
MP-DQN & 456 & $-$63237.50 & 201.00 & 23.6 & 78.00 & 0.040 & 242.30 \\
\bottomrule
\end{tabular}%
}
\end{table}

Every DQN/MP-DQN seed's switching frequency at $N=10$ (range
$0.030$-$0.096$) is lower than \emph{every} BMPP-DQN canonical seed's
(range $0.124$-$0.236$, Table~\ref{tab:oran-appendix-canonical}) --
not merely comparable, but consistently \emph{below} it. This is the
seed-level evidence behind Section~\ref{sec:oran-cadence-control}'s
central finding: BMPP-DQN's switching-frequency advantage over DQN/MP-DQN
at their own default (every-step) cadence is substantially a cadence
artifact, not an architectural one -- simple DQN/MP-DQN, once given
BMPP-DQN's own decision cadence, switch even less than BMPP-DQN itself.
